% Shared abstract and keywords. Every venue wrapper reads this file, so the
% four English editions of the paper cannot drift apart.
\newcommand{\capmabstract}{%
Vehicle cybersecurity research and regulation concentrate on the product:
ISO/SAE~21434 and UN~R155 govern the car, its ECUs and its back end. The
factory that builds the car is governed by nothing comparable, yet it is the
asset whose failure stops output, starves a continental supply chain and, in
the most expensive cyber event recorded in the United Kingdom, cost an economy
an estimated \pounds1.9 billion. This paper treats the automotive plant, rather
than the vehicle, as critical infrastructure, and asks a question that the
vehicle-centric literature cannot answer: by which paths does an attacker
actually reach the assembly line? We contribute (i)~a curated corpus of 20
publicly reported incidents in automotive manufacturing between 2017 and 2025,
each reconstructed as an asset-level path with an explicit confidence label;
(ii)~\capm, a cross-plane attack-path model that spans corporate IT, the
identity plane, cloud and SaaS, the industrial DMZ, MES/SCADA/PLC layers and
the supplier interface, instantiated as a reference architecture of 12 IEC
62443 zones, 48 assets and 91 attack steps labelled with 43 MITRE \attck{}
Enterprise and \attck{} for ICS techniques; (iii)~a probabilistic analysis
combining $k$-most-likely-path ranking with a Monte Carlo campaign simulation
that quantifies annualised loss; and (iv)~a control catalogue mapped to IEC
62443-3-3 and NIST CSF 2.0, evaluated by marginal value, necessity and
cost-effectiveness. Three results stand out. First, the path of least
resistance to a production stoppage never touches a controller: the most likely
route is $11.6\times$ more likely than the most likely route that reaches
Purdue level~2 or below, and the corpus agrees---55\% of incidents caused an
availability impact while only 10\% are reported to have interacted with
equipment below level~3. Second, recovery capability dominates prevention in
expected-loss terms: tested offline backups and a rehearsed degraded-production
mode account for the two largest single-control reductions, while SaaS token
governance and supplier assurance are the controls a mature programme can least
afford to remove. Third, partial hardening displaces rather than removes risk:
a supply-chain-and-cloud programme cuts espionage loss by 73\% but raises
ransomware loss by 26\% by pushing a rational actor onto the availability
objective. The model, corpus, experiments and figures are released as a
dependency-free open-source artefact.%
}

\newcommand{\capmkeywords}{%
critical infrastructure, automotive manufacturing, attack graphs, operational
technology, MITRE \attck{} for ICS, IEC 62443, supply-chain security, risk
quantification%
}
